%======================================================================
\section{The law: unreachable perfection forces usable entanglement}\label{sec:law}
%======================================================================

This section proves Theorem~A of the introduction and three further results: effective certificates for every fixed
target, undecidability of the dichotomy across games, and a finite protocol in a coherent white-box model. The certificates and
the protocol are proved for rational games in class~(S) and rational binary games in class~(C).

\subsection{The two classes}\label{sec:law-classes}

\begin{definition}[The classes of Theorem~A]\label{def:law-classes}
A finite game $G$, with a fixed question distribution $\mu$, belongs to class~(S) or class~(C) as follows.
\begin{enumerate}[label=(\Alph*),itemsep=2pt]
\item[(S)] \emph{Synchronous, with tested diagonals.} Both players have the same finite question and answer sets, unequal
answers lose on equal questions, and every diagonal pair $(x,x)$ has $\mu(x,x)>0$. The relations imposed are the vanishing
of $e_{xa}e_{yb}$ for the losing answer pairs $(a,b)$ on question pairs with $\mu(x,y)>0$ (the \emph{supported losing-pair
quotient}). A predicate merely called synchronous whose diagonals have zero weight is not in the class.
\item[(C)] \emph{Constraint-variable consistency.} There are finitely many variables with finite domains and finitely many
constraints. Alice receives a constraint and answers a satisfying assignment to its variables; Bob receives one variable of
that constraint and answers a value; they win if the values agree. Every constraint-variable incidence of every stated
constraint has positive probability. A variable that occurs in no constraint is never asked; we discard such variables,
which changes neither the correlations nor the demand, so every variable lies in some constraint. Native finite binary constraint-system games, including linear system games, are
the main examples.
\end{enumerate}
In both classes we assume that the finite-dimensional quantum value is one, so the set in~\eqref{eq:dG} is nonempty for
every $\eps>0$.
\end{definition}

For class~(C) let $\mathcal A(G)$ be the universal unital $C^*$-algebra generated by a projection-valued measure for each
variable, with the variable projections within each constraint commuting and every ordered product of projections
representing a forbidden assignment equal to zero. A finite-dimensional perfect strategy exists if and only if
$\mathcal A(G)$ has a nonzero finite-dimensional representation: given one, Bob measures the variable PVMs, Alice
measures the transpose of the joint PVM of the queried constraint, and the state is $\Phi_d$; the converse is
\cref{prop:normal-perfect} below in its finite-dimensional case. For class~(S) the same holds with the supported
losing-pair quotient, with Alice using transposes of the question PVMs.

\subsection{Theorem A}\label{sec:theoremA}

\begin{theorem}[Theorem A]\label{thm:A}
Let $G$ be a finite game in class~(S) or class~(C). Then
\[
G\text{ has no finite-dimensional perfect strategy}\quad\Longleftrightarrow\quad d_G^\to(\eps)\longrightarrow\infty\quad
(\eps\downarrow0).
\]
More precisely, if $G$ has no finite-dimensional perfect strategy, then for every integer $r\ge1$ and every $\gamma>0$ there
is $\eps_0(G,r,\gamma)>0$ such that every finite-dimensional mixed strategy of loss at most $\eps_0$ admits product local
isometries extracting $r$ Bell pairs with
fidelity at least $1-\gamma$; on any purification the output is within $\sqrt{2\gamma}$, in vector norm, of
$\Phi_{2^r}\otimes\chi$, all remaining systems, the environment included, being in the junk $\chi$. The isometries may
depend on the strategy and need not be efficient. The threshold is for fixed $r$ and $\gamma$; no rate in $r$ is asserted.
\end{theorem}

The direction ``$\Leftarrow$'' is immediate: a finite-dimensional perfect strategy is admissible at every $\eps$, and its
directional rates are at most the logarithm of its smaller local dimension. The same remark shows that the extraction clause
needs its hypothesis: a game with one question and one legal answer is won by a product state. The substance is
``$\Rightarrow$'', which we
prove in four steps: a tracial limit of Bob's moments, the exclusion of every finite type~I summand, bounded polynomial
approximation of Pauli matrices in the limit, and evaluation of those polynomials back on the original strategies.

\begin{proof}[Proof of \cref{thm:A}, class~(S)]
\emph{Step 1: actual mirrors and a tracial limit.} Let $(\rho_n)$ be finite-dimensional strategies of losses
$\eps_n\to0$. Dilate locally to projective measurements, and let $\psi_n$ purify the dilated state. For a question $x$,
$\sum_a\|(P^A_{xa}-P^B_{xa})\psi_n\|^2=2\Pr[a\ne b\mid(x,x)]\le2\eps_n/\mu(x,x)$, so every generator has an actual mirror
with a common error $\delta_n\to0$. Let $\mathcal A_0$ be the universal unital $C^*$-algebra of the free PVMs on the
finitely many questions, separable, and $\tau_n(F)=\langle\psi_n,F(P^B)\psi_n\rangle$. After passing to a subsequence,
$\tau_n\to\tau$ weak-$*$. Write $\theta$ for word reversal followed by substitution of Alice's PVMs; it extends to an
isometric $*$-antiautomorphism of $\mathcal A_0$ followed by a representation, and~\eqref{eq:opposite-poly} holds. For a
generator $e$ and a word $w$,
$\tau_n(ew)=\langle e^B\psi_n,w^B\psi_n\rangle$, and replacing $e^B\psi_n$ by $e^A\psi_n$, commuting $e^A$ past $w^B$ and
replacing it back costs at most $2\delta_n$. So $|\tau_n(ew)-\tau_n(we)|\le2\delta_n$, and by iteration, linearity and
density $\tau$ is a trace. If $(a,b)$ loses at $(x,y)$ with $\mu(x,y)>0$, then $\tau_n(e_{xa}e_{yb})$ is, up to $\delta_n$,
the actual probability of $(a,b)$ at $(x,y)$, which tends to zero; so $\tau(e_{xa}e_{yb})=0$ and, $\tau$ being a trace,
$\|e_{xa}e_{yb}\|_{2,\tau}^2=\tau(e_{xa}e_{yb})=0$. The tracial null space is a two-sided ideal, so in the GNS representation
$\pi_\tau$ every forbidden product vanishes as an operator. Put $\mathcal M=\pi_\tau(\mathcal A_0)''$, a finite von Neumann
algebra with faithful normal trace $t$ on a separable space.

\emph{Step 2: no finite type~I summand.} Suppose $\mathcal M$ had a nonzero type~I central summand. By the structure
theorem it contains a nonzero homogeneous piece $L^\infty(X,\nu)\otimes M_d(\C)$ for some finite $d$; there is no type
$\mathrm{I}_\infty$ piece in a finite algebra, and separability gives standard measure spaces. Represent the finitely many
generating projections on this piece by measurable fields of $d\times d$ matrices. The finitely many defining relations
(projections, PVM sums, forbidden products) hold almost everywhere, hence on one common conull set; at any point of it the
matrices are PVMs satisfying all relations, and with Alice's transposes and $\Phi_d$ they give a finite-dimensional
perfect strategy. This also excludes diffuse type~I centres. So, if $G$ has no finite-dimensional perfect strategy,
$\mathcal M$ is of type~II, not necessarily a factor, and for every $d$ it contains a unital copy of $M_d(\C)$: the
centre-valued trace divides $1$ into $d$ equivalent projections~\cite[Lemma~1.4]{CETW22}.

\emph{Step 3: bounded approximation.} Fix $r$, put $m=2r$, and let $q_1,\dots,q_m\in\mathcal M$ be the standard Pauli
reflections $X_1,\dots,X_r,Z_1,\dots,Z_r$ of a unital $M_{2^r}(\C)$, with $q_jq_k=s_{jk}q_kq_j$. Fix $\eta\in(0,1)$. By
Kaplansky's density theorem~\cite{Kap51} there are self-adjoint contractions $a_j\in\pi_\tau(\mathcal A_0)$ with
$\|a_j-q_j\|_{2,t}<\eta/2$. Each lifts to a self-adjoint contraction $h_j\in\mathcal A_0$ (lift, average with the adjoint,
and apply the clipping function $\max(-1,\min(t,1))$, which commutes with the quotient map); approximate $h_j$ in norm by a
self-adjoint $*$-polynomial $F_j$ within $\eta/2$. Then $\|F_j\|_{\mathcal A_0}\le2$, $\|\pi_\tau(F_j)-q_j\|_{2,t}<\eta$, and,
the trace making bounded multiplication continuous in $\|\cdot\|_{2,t}$,
\begin{equation}\label{eq:A-limit}
\|\pi_\tau(F_j^2-1)\|_{2,t}\le3\eta,\qquad\|\pi_\tau(F_jF_k-s_{jk}F_kF_j)\|_{2,t}\le6\eta .
\end{equation}

\emph{Step 4: back on the original strategies.} Fix these polynomials before $n$ grows, and put $H^B_j=F_j(P^B_n)$ and
$H^A_j=\theta(F_j)(P^A_n)$, self-adjoint of norm at most two on both parties. Moment convergence turns the squared vector
norms of the fixed defect polynomials on $\psi_n$ into their $\tau$-norms, so Bob's square and pair defects are eventually
at most $3\eta+o(1)$ and $6\eta+o(1)$. By~\eqref{eq:opposite-poly} the mirrors $\|(H^A_j-H^B_j)\psi_n\|$ tend to zero, and
since $\theta(F_jF_k-s_{jk}F_kF_j)=-s_{jk}(\theta F_j\theta F_k-s_{jk}\theta F_k\theta F_j)$ the same bounds hold for Alice.
Round locally, $Q^P_j=\sgn(H^P_j)$ with $\sgn(0)=1$; the scalar inequality $|\sgn t-t|\le|t^2-1|$ bounds the rounding error
on $\psi_n$ by the square defect. Inside a product, $\|(Q_j-H_j)H_k\psi_n\|\le\|H_k^{\mathrm{op}}(Q_j-H_j)\psi_n\|+\|Q_j-H_j\|
\|(H_k-H_k^{\mathrm{op}})\psi_n\|$, where $H_k^{\mathrm{op}}$ is the actual opposite-party mirror; so each product
replacement costs at most three times the rounding error plus three times the mirror error. Altogether the rounded
reflections satisfy, on both parties,
\[
\|(Q^P_jQ^P_k-s_{jk}Q^P_kQ^P_j)\psi_n\|\le\kappa_n,\quad\|(Q^A_j-Q^B_j)\psi_n\|\le\kappa_n,\quad
\limsup_n\kappa_n\le100\eta .
\]
\Cref{lem:fixed-size} now gives product local isometries whose output has singlet infidelity at most
$2r(2C_r+1)^2\kappa_n^2/4$ and, by~\eqref{eq:bell-hashing}, $\ED^{A\to B}(\rho_n)\ge r[1-2h(p)]$ with
$p=(2C_r+1)^2\kappa_n^2/4$. Letting $\eta\to0$ after $n\to\infty$, $\liminf_n\ED^{A\to B}(\rho_n)\ge r$, and for every
$\gamma$ the extraction has fidelity at least $1-\gamma$ for all large $n$.

\emph{Uniformity.} If for some finite $K$ there were strategies with losses tending to zero and $\ED^{A\to B}\le K$, apply
the above to that sequence with $r>K$: a contradiction. If for some $r,\gamma$ no threshold worked, there would be a
sequence with losses tending to zero none of which admits an $r$-pair extraction of fidelity $1-\gamma$: again a
contradiction. The vector-norm form follows from the fidelity form by projecting onto the Bell factor.
\end{proof}

\begin{proof}[Proof of \cref{thm:A}, class~(C)]
After local projective dilation let $B_{j,a}$ be Bob's variable-value projections and $C_{i,j,a}$ Alice's coarse-grained
projections of the joint assignment for constraint $i\ni j$. For fixed $i$ the $C_{i,j,a}$ commute and their forbidden
products vanish exactly, since Alice's answers are satisfying assignments. Positive probability of every incidence and
vanishing loss give $\|(C_{i,j,a}-B_{j,a})\psi_n\|\to0$. Choose once an incident constraint $i(j)$ for each variable and use
$A_{j,a}=C_{i(j),j,a}$ as the actual Alice mirrors of the universal Bob generators; distinct proxies need not commute.
Step~1 now runs on the universal free-PVM algebra of the variable measurements and gives a trace. A Bob word on the
variables of one constraint $i$ transfers to the reversed word in that constraint's own commuting $C_{i,j,a}$, so Bob's
within-constraint commutators and forbidden products have vector residuals tending to zero; their squared $\tau$-norms
vanish and the tracial GNS ideal makes them exact operator relations in $\mathcal M$. A finite-dimensional
representation of this constraint algebra would be a finite-dimensional perfect strategy, so Steps~2--4 apply verbatim.
\end{proof}

\begin{remark}[What is used, and what is not]\label{rem:A-scope}
The limit is a tracial GNS representation of finitely many approximate algebra elements, not an assumed spatial perfect
strategy and not a matrix ultraproduct with an unproved lift; every operation is evaluated back on the original
strategies, with their purifying environments. The exact tracial correspondence for perfect consistency is that of
Cleve--Mittal and of Paulsen and collaborators~\cite{CM14,PSSTW16,KPS18}; approximate tracial rounding of near-perfect
strategies is due to Vidick, Marrakchi--de la Salle, Paddock and Culf~\cite{Vid22,MdlS23,Pad22,Cul26}. Rounding to a
favorable tracial strategy is not by itself a lower bound on $\ED$ of an arbitrary original mixed state, and that is the
conclusion here. Approximating a strategy blockwise by maximally entangled ones would give a different state on each block, while
extraction needs one pair of local isometries on the original state, a difficulty discussed by Vernooij and
Zhao~\cite[Section~1.2]{VZ26}; the proof above avoids it, because the polynomials are evaluated on the original strategies
themselves. Nor does the proof use amenability of the limit trace: amenability need not pass from a trace on a universal algebra
to a quotient~\cite[Remark~5.15]{Zha24}, but the limit here is only a tracial state of the free algebra $\mathcal A_0$, the losing
products vanish in its GNS representation, and the polynomials are chosen in $\mathcal A_0$, where the moments converge. The
type dichotomy is the one behind the infinite-dimensional results compared in \cref{sec:intro-known}: Keyl, Matsui, Schlingemann
and Werner prove the infinite-one-copy-entanglement statement for every normal joint state of hyperfinite factors in Haag duality
that are not of type~I~\cite[Theorem~4.9]{KMSW06}, and van Luijk, Stottmeister, Werner and Wilming remove hyperfiniteness for
extraction from pure vector states~\cite[Theorem~54]{vLSWW25}; the mixed-state assertion for finite strategies is proved here
directly. No classification of arbitrary games is proved, and whether the conclusion extends beyond classes~(S) and~(C) is open
(\cref{sec:status}).
\end{remark}

\begin{remark}[Finite menus]\label{rem:finite-menu}
For fixed $(G,r,\gamma)$, finitely many polynomial tuples suffice for all sufficiently good strategies. The tracial
states of the game algebra with zero losing products form a compact set $T_G$; each $\tau\in T_G$ has a tuple whose
defects~\eqref{eq:A-limit} are below a fixed tolerance, which is a weak-$*$ open condition; compactness gives a finite
cover. If a sequence of strategies with losses tending to zero defeated every tuple of the cover, a moment subsequence
would converge into $T_G$ and into one of the open sets, a contradiction. The successful tuple may depend on the strategy.
\end{remark}

\subsection{Effective certificates}\label{sec:effective}

The compactness proof gives no threshold. For rational games, with binary variables in class~(C), a threshold can
nevertheless be computed, by a search that is guaranteed to halt under the promise.

\begin{theorem}[Effective fixed-target certificates]\label{thm:effective}
Let $G$ be a rational game in class~(S), or a rational finite binary constraint-variable game in class~(C), promised to
have finite-dimensional quantum value one and no finite-dimensional perfect strategy. For fixed $r$ and rational
$0<\gamma<1$, a partial algorithm halts and returns a rational $\eps_0>0$, a finite menu of rational universally
contractive self-adjoint polynomial arrays, each followed by local sign rounding and the fixed signed regular isometries
of \cref{lem:fixed-size}, and a finite rational tracial-moment infeasibility certificate, such that every strategy of loss
at most $\eps_0$ is decoded by some member of the menu with an $r$-pair output of fidelity at least $1-\gamma$. No useful
running time, closed-form dependence on $r$, or decision of the promise is asserted.
\end{theorem}

\begin{proof}
\emph{Universally contractive rational polynomials.} Enumerate rational word polynomials $q$ and put $p=\frac12(q+q^*)$,
$C=1+\sum_w(|\Re c_w|+|\Im c_w|)$, so $\|p\|\le C$ in every PVM representation. With $c(t)=\max(-1,\min(t,1))$ and the
Bernstein polynomial $b_{C,n}(t)=\sum_{k=0}^nc(-C+2Ck/n)\binom nku^k(1-u)^{n-k}$, $u=(t+C)/(2C)$, the polynomial
$F=b_{C,n}(p)$ has rational coefficients and satisfies $F=F^*$ and $\|F\|\le1$ exactly in every representation, because the
Bernstein weights are commuting positive operators summing to one. Also $\sup_{|t|\le C}|b_{C,n}(t)-c(t)|\le C/\sqrt n$, and
these $F$ are norm-dense in the self-adjoint unit ball of $\mathcal A_0$.

\emph{A finite cover.} Put $m=2r$, $C_r=m(1+2m)$ and the rational target $a=\gamma/[10^8(r+1)(2C_r+1)^2]$. For an array
$F=(F_1,\dots,F_m)$ let $d_F=\sum_j(F_j^2-1)^2+\sum_{j<k}(F_jF_k-s_{jk}F_kF_j)^*(F_jF_k-s_{jk}F_kF_j)$. By Steps~2--3 of the
proof of \cref{thm:A}, every perfect trace $\tau\in T_G$ has an enumerated array with $\tau(d_F)<a/2$, and compactness gives
a finite menu $F^{(1)},\dots,F^{(K)}$ with $T_G\cap\{\tau(d_{F^{(k)}})\ge a\ \forall k\}=\emptyset$.

\emph{A complete moment hierarchy, with two-sided PVM relations.} At level $N$ use moments $y_w$ for words of length at most
$2N$, with $y_1=1$, $y_{w^*}=\overline{y_w}$, a positive semidefinite moment matrix $[y_{u^*v}]_{|u|,|v|\le N}$, the relations
$y(u(e^2-e)v)=0$, $y(u(e-e^*)v)=0$ and $y(u(\sum_ae_{xa}-1)v)=0$ for \emph{all} left and right words $u,v$ within the degree
budget, cyclicity $y_{uv}=y_{vu}$ within the budget, $y(e_{xa}e_{yb})=0$ for supported losing pairs (for class~(C):
$y(Q^*Q)=0$ for the within-constraint commutators and forbidden products $Q$), boxes $|\Re y_w|,|\Im y_w|\le1$, and
$y(d_{F^{(k)}})\ge a$ for every $k$. This is a finite closed rational semialgebraic system. If it were feasible at every
level, a diagonal limit would be a positive normalized tracial functional; the boundedness of the generators in its GNS
construction uses $y(v^*(e^2-e)v)=0$, an instance of the two-sided relations, which gives $\|e[v]\|^2=\langle[v],e[v]\rangle$
and hence $\|e[v]\|\le\|[v]\|$. The limit would lie in $T_G$ and be bad for the whole menu, contradicting the cover. So some
finite level is infeasible, and by compactness it stays infeasible after relaxing only the trace and game equalities to
absolute error $\xi$, for some rational $\xi>0$. Exact real quantifier elimination decides each such instance~\cite{CM12}.

\emph{Transfer to near-winners.} Let $\mu_{\min}$ be the smallest positive incidence or diagonal weight; a strategy of
loss $\eps$ has generator mirrors $\delta\le\sqrt{2\eps/\mu_{\min}}$. Its Bob moments satisfy the PVM relations, positivity
and boxes exactly, the cyclic equalities up to $4N\delta$, and the game equalities up to $\eps/\mu_{\min}+\delta$ (for
class~(C), $\tau_B(Q^*Q)\le B_0^2\delta^2$ with $B_0=\max\{4,\max_i|i|\}$, by row transfer). Let $K_0\ge1$ be a rational bound on
$\sum_w|c_w||w|$ for every polynomial of the menu and its defects, and put
\[
z=\min\{1,a/K_0,\xi/(4N+1),\xi/2\},\qquad\eps_0=\min\{1,(\mu_{\min}/2)z^2,\mu_{\min}\xi/2\}
\]
(for class~(C), $z=\min\{1,a/K_0,\xi/(4N+1),\xi/(B_0+1)\}$ and $\eps_0=(\mu_{\min}/2)z^2$). For loss at most $\eps_0$ the
moments satisfy the relaxed system, so some array has $\tau_B(d_F)<a$: all Bob square and pair residuals are at most
$\sqrt a$, every raw polynomial mirror is at most $K_0\delta\le a$, and Alice's residuals are at most $2\sqrt a$. Sign
rounding and product replacement as in Step~4 give a common bound $\kappa=100\sqrt a$, and \cref{lem:fixed-size} gives
singlet infidelity at most $2r(2C_r+1)^2(100^2a)/4<\gamma$.

\emph{The search.} Enumerate arrays, menus, levels and rational $\xi\to0$, and run quantifier elimination on each finite
relaxed system; under the promise an infeasible candidate appears, and the algorithm outputs it with $\eps_0$. Every
returned certificate is sound regardless of the promise.
\end{proof}

Halting for one target does not decide the absence of finite-dimensional perfect strategies: a game with a large finite
perfect strategy can also satisfy a small-$r$ certificate. The rational Bernstein construction, tracial moment hierarchy,
quantifier elimination, type decomposition and Kaplansky density are established tools; the composed fixed-target
certificate for original mixed near-winners is what is new. The moment hierarchy is the tracial form of the semidefinite
hierarchies of Navascu\'es, Pironio and Ac\'{\i}n for quantum correlations and for noncommutative polynomial
optimization~\cite{NPA08,PNA10}; tracial moment matrices are studied by Burgdorf and Klep~\cite{BuKl12}. The relation
$y(u(e^2-e)v)=0$ must be imposed for all left and right words; one-sided imposition is not enough for the boundedness step.

\subsection{Undecidability of the dichotomy}\label{sec:undecidable}

\begin{theorem}[Bounded versus divergent demand is undecidable]\label{thm:undecidable}
There is a fixed bounded-output format, rational binary constraint-variable games with every incidence positively tested,
contexts of arity at most four, Alice at most sixteen answer labels and Bob two, such that, among games of this format
promised to have finite-dimensional quantum value one:
\begin{enumerate}[label=(\alph*),itemsep=1pt]
\item the games with bounded demand, $\lim_{\eps\to0}d_G^\to(\eps)<\infty$, are exactly those with a finite-dimensional
perfect strategy, and form a semidecidable set;
\item the games with divergent demand form an undecidable set that is not semidecidable.
\end{enumerate}
The numbers of questions may grow. The statement concerns finite game descriptions, not the distillability of a given
state or the capacity of a channel.
\end{theorem}

\begin{proof}
\emph{Inputs.} By Slofstra~\cite[Corollary~1.4]{Slo19}, it is undecidable whether a finite binary linear system game has a
perfect strategy in finite dimension. By~\cite[Theorem~1.1]{Slo19} and its construction~\cite[\S4,
\S5, Proposition~5.1 and Remark~5.5]{Slo19} there is a fixed linear system game $H$ with a perfect strategy in $C_{qa}$
but none in $C_{qs}$, hence with finite-dimensional quantum value one and no finite-dimensional perfect strategy. That
construction already has Alice eight answers and Bob two, with ternary equations, so no further arity reduction of $H$ is
needed.

\emph{Ternary inputs.} An input linear system $G$ is first replaced by one with contexts of at most three variables and
the same finite-dimensional operator solutions: for a row $x_1\cdots x_t=(-1)^b$ add, freshly for this row, $z_{ij}$ with
contexts $x_ix_jz_{ij}=1$ for $i<j$, prefix variables $p_2=z_{12}$ and $p_{k-1}x_kp_k=1$, and the unary row $p_t=(-1)^b$.
Given a solution put $z_{ij}=x_ix_j$ and $p_k=x_1\cdots x_k$; conversely the pair contexts force the $x_i$ to commute and
the prefixes recover the parity.

\emph{The selector.} With disjoint variable sets for $G$ and $H$ and one new flag $f$, let $K(G)$ have a context $C\cup\{f\}$
for each context $C$ of $G$, allowing exactly the assignments with $f=1$ or satisfying the constraint of $C$, and a context
$D\cup\{f\}$ for each context $D$ of $H$, allowing $f=0$ or the constraint of $D$; test every incidence, including every
occurrence of $f$, with positive rational probability. The flag occurs in every context, so in any operator solution its
projections $e_0,e_1$ commute with every variable and are central; all original context commutations remain required on
both branches.

\emph{Value one.} Play near-perfect finite strategies of $H$ with the flag deterministically $1$ and the $G$ variables set
to fixed scalars: $G$ contexts and all flag incidences pass exactly and the loss is a fixed weighted sum of the losses of
$H$. So every $K(G)$ has finite-dimensional value one, whatever $G$ is.

\emph{Perfection is preserved exactly.} If $G$ has a finite-dimensional operator solution, take $e_0=I$, keep it, and set
the $H$ variables to scalars. Conversely, a finite-dimensional perfect strategy of $K(G)$ gives a finite-dimensional
representation of its algebra (\cref{prop:normal-perfect}); split it by the central flag; the $e_1$ part would solve $H$
and is therefore zero, so $e_0=I$ and the $G$ operators solve $G$. Thus $K(G)$ has a finite-dimensional perfect strategy
if and only if $G$ does, and by \cref{thm:A} this happens if and only if $d_{K(G)}^\to$ stays bounded.

\emph{Undecidability and non-semidecidability.} A decision procedure for bounded demand on the promised class would decide
Slofstra's problem. Finite perfection is semidecidable: for $d=1,2,\dots$ decide by quantifier elimination whether a
$d$-dimensional operator solution exists. If divergence were also semidecidable, running both searches on $K(G)$ would
decide Slofstra's problem.

\emph{Bounded output.} The fallback $H$ is replaced by a fixed ternary $H_3$ by the same gadget; $H_3$ still has no finite
perfect strategy, and it has value one on the same states: for each auxiliary $z$ or $p$ representing a fixed product $w$
of $L$ original variables of one row, Bob uses $Q_w=\sgn(\frac12(U_w+U_w^*))$ with $U_w=\prod_{x\in w}B_x$, and Alice computes
the parity bits from her original row assignment; the estimates $\|(U_w-A_w)\psi\|,\|(U_w^*-A_w)\psi\|\le L\delta$,
$\|(F_w^2-I)\psi\|\le2L\delta$ and $\|(Q_w-A_w)\psi\|\le3L\delta$ show that every new incidence error tends to zero. With
$G_3$ and $H_3$ the selector has contexts of at most four variables including the common flag, so Alice's answers are
encoded in sixteen labels; invalid labels lose and can be relabelled to valid ones without increasing loss.
\end{proof}

The undecidable algebraic core is Slofstra's; the physical conclusion depends on \cref{thm:A}. The selector is an
elementary gated disjunction, and it is what keeps the promise: every $K(G)$ has value one through the fixed game $H_3$, whatever
$G$ is, and only the existence of a finite-dimensional perfect strategy depends on $G$. A closely related disjunction, in which the
players choose which of two games to play, is introduced by Man\v{c}inska and Schmidt: combining a game of value one without a
finite-dimensional perfect strategy, such as Slofstra's, with a synchronous game that self-tests a perfect finite-dimensional
strategy, they obtain an exact self-test that is not robust~\cite[Sections~4 and~5]{MaSc23}. Here the choice is a central flag
inside a Boolean constraint game, and the fixed nonattaining branch keeps the value-one promise for every input. A gapped problem,
value one against value at most $s<1$, would not serve here, since it does not separate attainment from nonattainment inside the value-one class.
For general two-player games such a gapped problem is undecidable by the theorem $\mathrm{MIP}^*=\mathrm{RE}$ of Ji,
Natarajan, Vidick, Wright and Yuen, which reduces the halting problem to deciding whether a game has entangled value one or at
most one half~\cite{JNVWY20}. Culf and
Mastel show that such gapped problems are already RE-complete for constraint-variable games of every Boolean constraint language
with NP-complete classical satisfiability, and for 3-colouring~\cite[Theorem~4.14 and Corollary~4.18]{CuMa25}. For example, their
3SAT instances have constraints of arity three and Boolean variable answers, and their 3-colouring instances have binary
constraints and a three-colour alphabet. These are value-gap results. Our dichotomy concerns finite perfect strategies versus
diverging usable entanglement under a value-one promise; we do not know whether its arity four and the sixteen Alice labels of
\cref{thm:undecidable} can be lowered. The theorem is compatible with \cref{thm:effective}: on a divergent instance every finite
target eventually has a certificate, but no finite prefix of such certificates recognizes divergence. We do not claim
that no computable precision modulus exists for a fixed game; a computably specified fixed value-one game has a
computable upper envelope of its demand by enumeration of finite witnesses.

\subsection{A finite protocol under coherent white-box access}\label{sec:protocol}

The effective certificate can be compiled into one finite protocol, independent of the hidden state and dimension, if
the players have coherent access to their own measurement circuits. This is a coherent white-box oracle model, \emph{not}
a device-independent one.

\begin{theorem}[White-box protocol]\label{thm:protocol}
Let $G$ be a rational game in class~(S) or a rational binary game in class~(C) satisfying the promise of
\cref{thm:effective}, and fix $r$ and rational $0<\gamma<1$. Assume: independent identically distributed copies of the
original state; fixed memoryless local measurement implementations; physical access to the payload and to reversible
Naimark circuits $V_x$ and $V_x^*$ with known outcome labels; pure ancillas; and coherent control of known circuits. A
partial compiler returns a loss threshold $\eps_*>0$ and \emph{one} finite protocol, independent of the hidden state, its
dimension and its matrix descriptions, with the following property. On every finite-dimensional mixed strategy of loss at
most $\eps_*$, Alice estimates her local defects on sacrificed copies, announces a menu index to Bob, and both apply
fixed coherent decoders to one retained copy, which yields $r$ qubits per side with singlet fraction at least $1-\gamma$;
all dirty work registers are kept in the accounting, and the guarantee holds against every purification of the retained
input. One sufficient number of sacrificed copies is
\begin{gather*}
K\Bigl\lceil2\Lambda_D^2b^{-2}\ln(2K/\zeta)\Bigr\rceil,\qquad b=\frac{\gamma}{10^8(r+1)(2C_r+1)^2},\\
\zeta=\frac\gamma8,\qquad m=2r,\qquad C_r=m(1+2m),
\end{gather*}
with $K$ the menu size and $\Lambda_D$ the recorded coefficient bound of the defect polynomials. All communication is from
Alice to Bob. The implementation uses a linear combination of unitaries, a two-reflection walk with phase-estimation
sign approximation, and the reversible measurement access; the gate costs of the unknown oracles are not bounded.
\end{theorem}

The proof is in \cref{app:protocol}. The theorem does not certify near-winning behaviour autonomously: local tests
cannot distinguish an entangled input from a product input with the same marginal. Nor does it leave the environment's
marginal unchanged after conditioning on the selection transcript; only the retained copy's $AB$ marginal is unchanged,
and the bound holds for every conditional purification. A device that offers only classical outputs, or only destructive
measurements, does not satisfy the interface: from classical outputs and fresh local ancillas only separable states can be
prepared.

Universal distillation from unknown mixed states is known in a fixed finite dimension. Takagi, Watanabe, Matsuura, Arai and
Hayashi~\cite{TWMAH26} achieve the distillable entanglement of an arbitrary unknown state through tomography, state merging for
compound sources and uniform convergence, with almost one-way communication: Bob sends back tomographic information at a
vanishing rate~\cite[Theorem~8 and Appendix~B.3(f)]{TWMAH26}. Tirone, Mele, Giovannetti and Lami~\cite{TMGL26} achieve it
independently with tomography by local operations and classical communication and a robust adaptive distillation routine. Both
build on protocols for compound sources~\cite{BJ14} and the earlier work cited there. So the possibility of distilling from an
unknown state is not what \cref{thm:protocol} adds. It instead uses a fixed-game promise and reversible coherent measurement
access to compile a finite dimension-uniform menu, selects a decoder with Alice-local data, and extracts a fixed target from one
retained input copy, using only communication from Alice to Bob. It does not optimize a universal asymptotic rate and does not
provide an efficient device-independent procedure.
